\section{Finite generation of rational log canonical rings}
\label{sec:canonical-rings}

The preceding results also settle finite generation for rational lc pairs.
The finite-generation conjecture asks whether finitely many rounded
pluricanonical sections generate the log canonical ring
\cite[Conjecture~A]{FujinoGongyoRings2014}. The assertion below concerns
the full rounded ring, not merely a sufficiently divisible subring, and
requires no nefness assumption.
For the projective absolute statement, we combine the ordinary minimal-model
theorem with log abundance and keep track of every rounded degree. For the
proper relative statement over \(\C\), we apply Fujino--Gongyo's published
reduction from good minimal models.
In that relative statement the morphism is only required to be proper;
the source need not be projective over \(\C\).

For an integral Weil divisor \(G\) on a normal variety \(X\),
\(\OO_X(G)\) denotes its divisorial sheaf. For a rational divisor \(D\)
on \(X\) and nonnegative integers \(m,n\), the inclusions
\(\lfloor mD\rfloor+\lfloor nD\rfloor\leq\lfloor(m+n)D\rfloor\)
give the usual multiplication on the section algebras below.

\begin{samepage}
\begin{corollary}[Finite generation for rational lc pairs]
\label{cor:lc-finite-generation}
\leavevmode
\begin{enumerate}[label=\textup{(\roman*)}]
\item Let \(k\) be an algebraically closed field of characteristic zero,
let \((X,\Delta)\) be a normal projective log canonical pair over \(k\), and assume
that \(\Delta\geq0\) is rational and \(D=K_X+\Delta\) is \(\Q\)-Cartier.
Then
\[
 R(X,D)=\bigoplus_{m\geq0}
 H^0\bigl(X,\OO_X(\lfloor mD\rfloor)\bigr)
\]
is a finitely generated \(k\)-algebra.

\item Let \((X,\Delta)\) be a log canonical pair over \(\C\), with
\(\Delta\geq0\) rational and \(K_X+\Delta\) \(\Q\)-Cartier. For every
proper morphism \(f:X\to S\) onto an algebraic variety \(S\), the relative
log canonical algebra
\[
 \mathcal R(X/S,K_X+\Delta)
 =\bigoplus_{m\geq0}f_*\OO_X\bigl(\lfloor m(K_X+\Delta)\rfloor\bigr)
\]
is a finitely generated \(\OO_S\)-algebra.
\end{enumerate}
\end{corollary}
\end{samepage}

\begin{proof}
In dimension zero, \(X\) is a point, and in \textup{(ii)} so is \(S\).
The boundary and canonical divisor are zero, and the two algebras are
\(k[t]\) and \(\OO_S[t]\), respectively, with \(\deg t=1\).
We therefore assume \(\dim X\geq1\).

We first prove \textup{(i)}. If \(D\) is not pseudo-effective, a nonzero
section \(s\) in degree \(m>0\) would give
\[
 \operatorname{div}(s)+mD
 =\bigl(\operatorname{div}(s)+\lfloor mD\rfloor\bigr)
   +\bigl(mD-\lfloor mD\rfloor\bigr)\geq0.
\]
This would make \(D\) pseudo-effective. Thus all positive graded pieces
vanish in this case. Since \(H^0(X,\OO_X)=k\), we have \(R(X,D)=k\).

Suppose now that \(D\) is pseudo-effective. The ordinary-pair conclusion
of the companion minimal-model theorem
\cite[Corollary~1.2]{OpenAIMinimalModels2026} gives a normal projective lc
pair \((Y,\Delta_Y)\) and a birational map \(\phi:X\dashrightarrow Y\)
extracting no divisors, with \(\Delta_Y=\phi_*\Delta\) and
\(D_Y=K_Y+\Delta_Y\) \(\Q\)-Cartier and nef. Log discrepancies do not decrease.
Theorem~\ref{thm:main} makes the rational adjoint \(D_Y\) semiample.

On a common smooth resolution \(p:W\to X\), \(q:W\to Y\), choose
compatible canonical divisors. Then
\begin{equation}\label{eq:ring-comparison}
 p^*D=q^*D_Y+E,\qquad E\geq0,\qquad E\text{ is }q\text{-exceptional}.
\end{equation}
Indeed, the coefficient of \(E\) at a prime divisor \(F\) is
\(a(F;Y,\Delta_Y)-a(F;X,\Delta)\), which is nonnegative by the discrepancy
comparison. Suppose \(F\) is not \(q\)-exceptional. Since \(\phi\) extracts no
divisors, \(F\) is the strict transform of a divisor on \(X\). The
pushed-forward boundary gives equality of these discrepancies. This proves
the asserted support.

For any normal variety \(Z\), rational Cartier divisor \(G\), and
\(m\geq0\), the nonzero global sections of \(\OO_Z(\lfloor mG\rfloor)\) are
exactly the rational functions \(s\) satisfying
\[
 \operatorname{div}_Z(s)+mG\geq0.
\]
This is equivalent to the usual inequality with \(\lfloor mG\rfloor\)
because the coefficients of \(\operatorname{div}_Z(s)\) are integers.
If the inequality holds on \(X\), pull back its effective \(\Q\)-Cartier
left side to \(W\) and push forward by \(q\). Equation~\eqref{eq:ring-comparison}
and \(q_*E=0\) give the inequality on \(Y\). Conversely, pull back the
inequality on \(Y\), add \(mE\), and push forward by \(p\). Thus, inside
the common function field,
\[
 H^0\bigl(X,\OO_X(\lfloor mD\rfloor)\bigr)
 =H^0\bigl(Y,\OO_Y(\lfloor mD_Y\rfloor)\bigr)
 \qquad(m\geq0).
\]
These identifications preserve multiplication, so they identify the full
graded rings.

It remains to use semiampleness on \(Y\). Choose \(r>0\) such that
\(rD_Y\) is Cartier and globally generated. Its complete linear system
defines \(\psi:Y\to\PP^N_k\) with
\(\OO_Y(rD_Y)\simeq\psi^*\OO_{\PP^N}(1)\). Let
\(A=k[T_0,\ldots,T_N]\) act on \(R(Y,D_Y)\) by pulling back the coordinate
sections, each in degree \(r\). For \(0\leq i<r\), put
\(\mathcal F_i=\psi_*\OO_Y(\lfloor iD_Y\rfloor)\), which is coherent on
\(\PP^N_k\) because \(\psi\) is proper. Since
\[
 \lfloor(\ell r+i)D_Y\rfloor
 =\ell(rD_Y)+\lfloor iD_Y\rfloor\qquad(\ell\geq0),
\]
projection formula identifies the part in degrees congruent to \(i\)
modulo \(r\), as an \(A\)-module, with
\[
 \bigoplus_{\ell\geq0}H^0\bigl(\PP^N_k,\mathcal F_i(\ell)\bigr).
\]
This is a finite \(A\)-module by the standard finite-generation theorem
for section modules \cite[Lemma~30.14.1(5)]{StacksSectionModules}.
There are only \(r\) residue classes, so \(R(Y,D_Y)\) is finite as an
\(A\)-module and hence finitely generated as a \(k\)-algebra. This proves
\textup{(i)}.

For \textup{(ii)}, set \(n=\dim X\). In the notation of Fujino--Gongyo,
Conjecture \(C_{\leq n-1}\) asks for good minimal models of projective
\(\Q\)-factorial dlt pairs with effective real boundary and pseudo-effective
adjoint in dimensions at most \(n-1\). Theorem~\ref{thm:good-models-complex}
supplies this statement. Its good-model convention is the one fixed in
Section~\ref{sec:induction}, with a \(\Q\)-factorial dlt model as in
Proposition~\ref{prop:inductive-reduction}. Theorem~1.1 of
Fujino--Gongyo \cite{FujinoGongyoRings2014} therefore gives their
Conjecture \(B_n\): finite generation of the log canonical ring for
projective rational plt pairs of dimension \(n\) whose adjoint is big
and whose boundary has irreducible round-down. Their relative reduction,
Corollary~1.4, applies to the given rational lc pair and proper morphism
onto \(S\), and gives exactly \textup{(ii)} without requiring \(X\) to
be projective over \(\C\).
Their Theorem~1.1 also gives \textup{(i)} directly when \(k=\C\).
\end{proof}

The good-model input allows real boundaries, but both finite-generation
statements retain rational boundaries. In the absolute proof, rationality
is what leaves finitely many residue-class section modules. The relative
assertion is finite generation over the algebraic base \(S\) in the complex
setting; it is not a claim about rounded rings for real boundaries or about
analytic bases. No uniform degree bound for generators is asserted.
